\chapter*{Progression Annuelle \& Référentiel}
\addcontentsline{toc}{chapter}{Progression Annuelle \& Référentiel}
\markboth{PROGRESSION ANNUELLE}{RÉFÉRENTIEL 3PM}

\begin{cadreretenu}[Objectifs de la Classe de 3\textsuperscript{ème} Prépa-Métiers]
La classe de 3\textsuperscript{ème} Prépa-Métiers permet aux élèves de consolider le socle commun de connaissances, de compétences et de culture tout en découvrant des environnements professionnels variés. En mathématiques, l'accent est mis sur :
\begin{itemize}
    \item La consolidation des automatismes de calcul et du raisonnement logique.
    \item La résolution de situations concrètes issues des métiers (BTP, Industrie, Tertiaire, Artisanat).
    \item L'utilisation efficace des outils numériques (Calculatrice, Tableur, GeoGebra, Scratch).
    \item La préparation rigoureuse à l'épreuve de Mathématiques du \textbf{Diplôme National du Brevet (DNB Pro)}.
\end{itemize}
\end{cadreretenu}

\vspace{0.5cm}

\section*{Grille Nationale des Compétences Évaluées}

\noindent
\begin{tabularx}{\linewidth}{|c|l|X|}
\hline
\rowcolor{slate800} \textcolor{white}{\textbf{Code}} & \textcolor{white}{\textbf{Compétence}} & \textcolor{white}{\textbf{Description synthétique}} \\
\hline
\rowcolor{slate50} \capSApproprier & \textbf{S'approprier} & Extraire l'information utile, reformuler une situation, identifier les données d'un problème métier. \\
\hline
\rowcolor{white} \capAnalyser & \textbf{Analyser / Raisonner} & Choisir un modèle (linéaire, géométrique, statistique), émettre une conjecture, élaborer une démarche. \\
\hline
\rowcolor{slate50} \capRealiser & \textbf{Réalisér} & Effectuer un calcul numérique ou algébrique, utiliser la calculatrice, construire une figure, exécuter un algorithme. \\
\hline
\rowcolor{white} \capValider & \textbf{Valider} & Contrôler la vraisemblance d'un résultat, vérifier des unités de mesure, critiquer une hypothèse. \\
\hline
\rowcolor{slate50} \capCommuniquer & \textbf{Communiquer} & Rédiger une solution claire, argumenter à l'oral ou à l'écrit, expliciter les étapes de son raisonnement. \\
\hline
\end{tabularx}

\vspace{0.8cm}

\section*{Tableau de Progression Annuelle (5 Périodes)}

\noindent
\begin{tabularx}{\linewidth}{|c|l|X|c|}
\hline
\rowcolor{colPrimary} \textcolor{white}{\textbf{Période}} & \textcolor{white}{\textbf{Domaine}} & \textcolor{white}{\textbf{Chapitres \& Notions au Programme}} & \textcolor{white}{\textbf{Éval.}} \\
\hline
\rowcolor{colDefBg} \textbf{Période 1} & \textbf{Nombres \& Calculs} & $\bullet$ Divisibilité, décomposition en facteurs premiers, fractions irréductibles.\newline $\bullet$ Puissances de 10, écriture scientifique et ordres de grandeur.\newline $\bullet$ \textbf{Calcul Littéral} : conventions, développement $k(a+b)$, $(a+b)(c+d)$, réduction. & DS 1 \\
\hline
\rowcolor{white} \textbf{Période 2} & \textbf{Algèbre \& Géométrie} & $\bullet$ \textbf{Équations du 1\textsuperscript{er} degré} : résolution algébrique $ax+b=c$, problèmes concrets.\newline $\bullet$ Proportionnalité, pourcentages d'évolution, ratios et échelles.\newline $\bullet$ \textbf{Théorème de Pythagore} : calcul de longueurs, réciproque et contraposée. & DS 2 \\
\hline
\rowcolor{colDefBg} \textbf{Période 3} & \textbf{Fonctions \& Espace} & $\bullet$ \textbf{Fonctions affines et linéaires} : $f(x)=ax+b$, représentations graphiques, lecture de coefficients.\newline $\bullet$ \textbf{Trigonométrie} dans le triangle rectangle ($\cos, \sin, \tan$).\newline $\bullet$ Repérage dans le plan et dans l'espace (coordonnées 2D/3D). & Brevet Blanc 1 \\
\hline
\rowcolor{white} \textbf{Période 4} & \textbf{Géométrie \& Données} & $\bullet$ \textbf{Théorème de Thalès} \& Agrandissement-Réduction (rapports $k, k^2, k^3$).\newline $\bullet$ \textbf{Statistiques à 1 variable} : moyenne pondérée, médiane, étendue, fréquences.\newline $\bullet$ \textbf{Probabilités} : expériences aléatoires simples et à 2 épreuves, arbres. & DS 3 \\
\hline
\rowcolor{colDefBg} \textbf{Période 5} & \textbf{Grandeurs \& Synthèse} & $\bullet$ \textbf{Aires, Volumes \& Conversions} : prismes, cylindres, cônes, pyramides, sphères.\newline $\bullet$ \textbf{Algorithmique \& Scratch} : boucles, conditions, variables, déplacements.\newline $\bullet$ \textbf{Entraînement intensif aux annales du DNB Pro}. & DNB PRO \\
\hline
\end{tabularx}

\newpage
